\section{Experimental Setup}
\label{sec:experimental_setup}

\subsection{Training data and model variants}

The training set is derived from trajectories generated by the nonlinear trajectory optimizer on the target Oval-track scenarios. These trajectories provide expert demonstrations under the same vehicle model and constraint structure used downstream. The final dataset combines four shards: 31,710 nominal solver trajectories, 400 hard-repair trajectories, 1,000 post-projection repair trajectories, and 300 failure-conditioned repair trajectories. The repair shards were introduced to expose the model to rollout states that do not appear in the nominal expert distribution.

The main model variants follow a progression from nominal offline training to increasingly deployment-aware fine-tuning. Earlier runs train on nominal data alone. Later runs incorporate larger context windows and targeted repair shards, including post-projection and failure-conditioned examples. The purpose of this progression is not to compare many unrelated models, but to test a specific hypothesis: whether better offline trajectory modeling, especially on rollout-relevant failure modes, leads to improved downstream solver utility.

\subsection{Training configuration}

Training is performed offline on fixed datasets of solver-generated trajectories using the architecture described in Section~2.2. The training pipeline is implemented in PyTorch, optimized with AdamW, logged through TensorBoard/JSON artifacts, and configured for CUDA execution with mixed precision in the final GPU training scripts. Validation loss is tracked throughout training and later used only as an intermediate model-quality metric rather than as the primary success criterion. The main progression of training interventions is summarized in Table~\ref{tab:training_summary}.

\begin{table}[t]
\centering
\caption{Representative training progression across the main DT variants. Validation loss is the best logged value for each checkpoint.}
\label{tab:training_summary}
\setlength{\tabcolsep}{4pt}
\small
\resizebox{\linewidth}{!}{%
\begin{tabular}{lccc}
\toprule
Model / Run & Data Regime & Val. Loss & Note \\
\midrule
Larger-context parent & nominal + hard-repair & 0.002835 & Best parent checkpoint \\
Post-projection fine-tune & + post-proj. repair & 0.003170 & Best deployment setting before failure FT \\
Lateral-stability fine-tune & + lateral weighting & 0.002773 & Improved offline stability \\
Failure-conditioned fine-tune & + failure-conditioned repair & 0.002364 & Best offline checkpoint \\
\bottomrule
\end{tabular}
\vspace{-1ex}
}
\end{table}

\subsection{Deployment pipeline and metrics}

Each model is evaluated as a warm-start generator for the downstream nonlinear solver. In the final project phase, the optimizer used for expert generation and warm-start evaluation is FATROP. For a given benchmark instance, the DT produces an autoregressive candidate trajectory, which is validated by the wrapper, optionally projected, and either accepted or rejected. Accepted trajectories are submitted to FATROP as warm-starts; rejected trajectories trigger fallback to the baseline initializer.

The primary deployment metrics are chosen to measure solver utility rather than imitation quality. These include:
\begin{itemize}
    \item warm-start acceptance rate,
    \item solver success rate,
    \item solver iterations,
    \item solver wall-clock time,
    \item total pipeline time including DT inference.
\end{itemize}
Additional diagnostic metrics quantify projection burden and rollout validity, helping distinguish between trajectories that are merely accepted and those that materially improve the downstream solve.

Offline validation loss is also reported, but only as an intermediate metric. The central question is whether lower validation loss corresponds to lower downstream optimization cost.

\subsection{Benchmarks and evaluation scope}

The evaluation uses both obstacle-free and obstacle scenarios on the target racing track. The main controlled benchmark consists of corrected 10-scenario gates for the larger-context parent and post-projection fine-tune. A later smoke benchmark evaluates the failure-conditioned model on 3 obstacle-free and 3 obstacle scenarios after the final debugging pass. These experiments are intentionally narrow in scope: they focus on a fixed vehicle model, a single track family, and a limited set of benchmark scenarios chosen to test whether DT warm-starts provide a practical advantage over the baseline initializer in this setting. The main benchmark constants are summarized in Table~\ref{tab:experiment_params}.

\begin{table}[h]
\centering
\caption{Main experiment parameters for the final Oval-track benchmark. Vehicle parameters come from the GTI model configuration used throughout the project.}
\label{tab:experiment_params}
\setlength{\tabcolsep}{4pt}
\small
\resizebox{\linewidth}{!}{%
\begin{tabular}{ll}
\toprule
Parameter & Value \\
\midrule
Track & Oval\_Track\_260m \\
Track length & 260 m \\
Track width & 6.0 m \\
Planning horizon & $N=120$ nodes \\
Vehicle & Volkswagen Golf GTI single-track model \\
Vehicle mass & 1868 kg \\
Wheelbase & 2.63 m \\
Max steering angle & 27 deg \\
Longitudinal force bounds & $[-15,\ 10]$ kN \\
Road condition & Dry asphalt \\
Nominal tire friction & $\mu_f=0.87,\ \mu_r=1.03$ \\
Obstacle scenarios & 0 or 1--4 circular obstacles \\
Obstacle radius range & 0.8--1.5 m \\
Obstacle margin & 0.3 m \\
Obstacle clearance buffer & 0.3 m \\
Solver & FATROP \\
\bottomrule
\end{tabular}
}
\end{table}
